% Διάμεσο μήκος λύσης ανά βάθος ανακατέματος και επιλυτή. Περιεχόμενο
% σχήματος· το περιβάλλον figure και η λεζάντα ορίζονται στο κεφάλαιο 9.
% Δεδομένα: figures/solution_length_depth_data.tex. Η σειρά Korf
% τερματίζει στο βάθος 10 (βαθύτερες περιπτώσεις χωρίς γραμμή διαμέσου).
\begin{tikzpicture}
\begin{axis}[
  width=0.85\textwidth,
  height=0.62\textwidth,
  xlabel={Βάθος ανακατέματος},
  ylabel={Διάμεσο μήκος λύσης},
  xmin=-0.5, xmax=20.5,
  xtick={0,5,10,15,20},
  minor x tick num=4,
  ymin=0, ymax=40,
  ytick={0,5,10,15,20,25,30,35,40},
  grid=major,
  grid style={dotted, gray!50},
  legend pos=north west,
  legend cell align=left,
  legend style={font=\footnotesize},
]
% Διαγώνιος αναφοράς y = x (βέλτιστο μήκος ίσο με το ονομαστικό βάθος).
\addplot[forget plot, gray, thick] coordinates {(0,0) (20,20)}
  node[pos=0.85, sloped, below, font=\scriptsize, text=gray] {βέλτιστο};
\addplot[blue, thick, mark=*, mark size=1.6pt] coordinates {
  (0,0) (1,1) (2,2) (3,3) (4,4) (5,5) (6,4) (7,7) (8,8) (10,18) (15,22) (20,23)
};
\addlegendentry{Kociemba (native)}
\addplot[red, thick, mark=square*, mark size=1.6pt] coordinates {
  (0,0) (1,1) (2,2) (3,3) (4,4) (5,5) (6,4) (7,7) (8,8) (10,9) (15,21) (20,21)
};
\addlegendentry{Kociemba (adapter)}
\addplot[green!55!black, thick, mark=triangle*, mark size=2pt] coordinates {
  (0,0) (1,1) (2,2) (3,3) (4,4) (5,5) (6,4) (7,7) (8,8) (10,9)
};
\addlegendentry{Korf}
\addplot[orange!90!black, thick, mark=diamond*, mark size=2pt] coordinates {
  (0,0) (1,1) (2,2) (3,3) (4,5) (5,17) (6,4) (7,25) (8,32) (10,31) (15,32) (20,26)
};
\addlegendentry{Thistlethwaite}
\end{axis}
\end{tikzpicture}
